\chapter{The terms equation (1) leaves out}

This chapter shows that the equation Fefferman poses~\cite{src:Fefferman} is not a reduction of
the balance laws of a real fluid but a different equation. It writes out every term the passage to
it removes, proves that the removed terms are not zero, and shows what putting the fluid in a
periodic box drops. It rests on equation (1) and condition (2) of Fefferman's statement, the
continuum balance laws of mass, momentum and internal energy, the hypotheses of Theorem~1.1 of the
Navier--Stokes paper~\cite{src:OpenAI-2026-Navier-Stokes} as the previous chapter quotes them, and
the first five pages of Lienstromberg, Schiffer and
Schubert~\cite{src:Lienstromberg-Schiffer-Schubert-2025}.

\section{The question}

Fefferman writes equation (1) in components,
\[
\frac{\partial u_i}{\partial t} + \sum_{j=1}^{n} u_j \frac{\partial u_i}{\partial x_j}
= \nu \Delta u_i - \frac{\partial p}{\partial x_i} + f_i(x,t),
\]
with condition (2), $\nabla\cdot u = 0$. It carries one constant, $\nu$. It carries no density and
no temperature. The fluid is told its viscosity and is never asked for it.

A term may be dropped from an equation in two cases and only two: when it is zero, or when it is
bounded by something small and the bound travels with the equation. This chapter writes the
removed terms out in vector form, where each one is a single product that cannot be overlooked,
and proves for each one that neither case holds for a fluid whose viscosity or density depends on
its temperature. Water is such a fluid.

The fluid is taken to be a continuum at every scale, however small. So is Fefferman's:
his $u$ is smooth on $\R^3\times[0,\infty)$. Nothing below departs from that. The fluid is also
taken to fill all of $\R^3$ unless a container has been measured. Nothing is known about a
container before its effect on the motion is read, and the last proved section shows what putting
the fluid in a box costs.

This is the first observation of the thought experiment behind it. Every term is written in vector
form, and nothing is held fixed that the fluid has not been asked for. The observations that follow
from it go in later chapters. The next one, \emph{What the average keeps}, places the molecules
in this field and bounds where an average of the two over a cell holds.

\section{The balance laws, kept whole}

Write $\rho$ for density, $u$ for velocity, $T$ for temperature, $\theta = \nabla\cdot u$ for the
rate of volume change, and $D = \tfrac12(\nabla u + \nabla u^{\mathsf T})$ for the rate of strain.
The continuum balance of mass and of momentum are
\[
\partial_t\rho + \nabla\cdot(\rho u) = 0,
\qquad
\rho\,(\partial_t u + (u\cdot\nabla)u) = \nabla\cdot\sigma + f,
\]
where $\sigma$ is the stress. The Newtonian stress law is the stress linear and isotropic in $D$,
\[
\sigma = -p\,I + \lambda\,\theta\,I + 2\mu D,
\]
with $\mu$ the shear viscosity and $\lambda$ the second coefficient. Neither coefficient is a
constant of nature. Each is a measured function of the state of the fluid, and the one this
chapter uses is $\mu = \mu(T)$.

Where $\theta = 0$, the balance of internal energy is
\[
\rho c\,(\partial_t T + u\cdot\nabla T) = \nabla\cdot(k\nabla T) + 2\mu\,|D|^2,
\]
with $c$ the specific heat, $k$ the thermal conductivity, and $|D|^2 = \sum_{i,j} D_{ij}^2$. The
last term is the heat that shear puts into the fluid, and it is never negative. These three balances
are the standard ones and are not derived here.

\section{The identity}

\textbf{Proposition 1.} For smooth $u$, $p$, $\mu$ and $\lambda$,
\[
\nabla\cdot\sigma = -\nabla p + \mu\,\Delta u + (\lambda+\mu)\,\nabla\theta + \theta\,\nabla\lambda
+ 2D\,\nabla\mu .
\]

\emph{Proof.} In components,
$\partial_j\sigma_{ij} = -\partial_i p + \partial_i(\lambda\theta) + \partial_j(2\mu D_{ij})$.
The second term is $\lambda\,\partial_i\theta + \theta\,\partial_i\lambda$. The third is
$2D_{ij}\,\partial_j\mu + 2\mu\,\partial_j D_{ij}$, and
$2\partial_j D_{ij} = \Delta u_i + \partial_i\theta$. Collecting gives the statement.
\hfill\qedsymbol

Set the density to one, $\theta$ to zero and $\nabla\mu$ to zero, and the momentum balance with the
Newtonian stress law becomes equation (1) with $\nu = \mu$. The three terms that leave are
\[
R = (\lambda+\mu)\,\nabla\theta + \theta\,\nabla\lambda + 2D\,\nabla\mu ,
\]
and equation (1) is exact at a point if and only if $R$ is zero there.

In components the shear part alone of $\partial_j\sigma_{ij}$ expands to twelve terms for each $i$,
and the three that carry $\partial_j\mu$ sit among nine that do not. In vector form the same three
are one product, $2D\nabla\mu$. The notation does not make a term large. It makes a term impossible
to drop without writing down the condition that drops it.

\section{What each condition asserts}

\begin{center}
\begin{tabular}{@{}p{0.24\textwidth}p{0.30\textwidth}p{0.38\textwidth}@{}}
\toprule
Condition & Term it removes & What it asserts about the fluid \\
\midrule
$\theta = 0$, condition (2) & $(\lambda+\mu)\nabla\theta + \theta\nabla\lambda$
& no element of the fluid changes volume \\[0.5em]
$\nabla\mu = 0$ & $2D\nabla\mu$ & viscosity is the same at every point \\[0.5em]
$\rho \equiv 1$ & the variation of $\rho$ in the inertia
& density is the same at every point and time \\[0.5em]
no energy balance & every coupling to $T$ & temperature does not enter the motion \\[0.5em]
the stress law itself & any dependence of $\mu$ on $|D|$ & stress is linear in the rate of strain \\
\bottomrule
\end{tabular}
\end{center}

Condition (2) removes $\lambda$ entirely. Whether the second coefficient is zero is therefore not
answered by equation (1). It is removed from the question.

\section{A viscous flow makes its own temperature gradient}

\textbf{Proposition 2.} Let $u$, $T$ be smooth on $\R^3\times I$ for an interval of time $I$,
satisfy the energy balance above with $\theta = 0$, $\rho c > 0$ and $\mu > 0$, and let $D(x,t)\to 0$
as $|x|\to\infty$ for each $t$. If $\nabla T = 0$ everywhere on $\R^3\times I$, then $u = 0$ on
$\R^3\times I$.

\emph{Proof.} With $\nabla T = 0$ the temperature is a function of time alone, the conduction term
vanishes, and the balance reads $\rho c\,T'(t) = 2\mu\,|D(x,t)|^2$ for every $x$. The left side does
not depend on $x$. Letting $|x|\to\infty$ sends the right side to zero. Then $T'(t) = 0$, and
$|D| = 0$ at every point. A smooth field with $D = 0$ is a rigid motion, $u = a(t) + W(t)x$ with
$W^{\mathsf T} = -W$. Decay at infinity forces $a = 0$ and $W = 0$. \hfill\qedsymbol

Any motion of a viscous fluid that decays at infinity, motion with compact support among them,
produces a temperature gradient within every interval of time it moves. The fluid writes its own
temperature field. Equation (1) has nowhere to put it.

\section{Dropping the gradient of viscosity}

\textbf{Proposition 3.} Let $\mu = \mu(T)$ and $\theta = 0$. The term removed by setting
$\nabla\mu = 0$ is $2\mu'(T)\,D\,\nabla T$. It is zero at a point if and only if $\mu'(T) = 0$ there
or $\nabla T$ lies in the null space of $D$. At a point where $\det D \neq 0$ and $\mu'(T)\neq 0$, it
is zero if and only if $\nabla T = 0$.

\emph{Proof.} The chain rule gives $\nabla\mu = \mu'(T)\nabla T$. A matrix with nonzero determinant
has trivial null space. \hfill\qedsymbol

For water the factor $\mu'(T)$ is not small. IAPWS R12-08~\cite{src:IAPWS-R12-08}, the formulation for the viscosity of
ordinary water, evaluated by a program that reproduces its verification table to every
printed digit, gives $1.0016$ mPa\,s at $20\,^\circ$C and atmospheric pressure and
$\mu'/\mu = -0.0245$ per kelvin there: about two and a half percent per degree.

The term was dropped for a stated reason. Stokes~\cite{src:Stokes-1845}, in his Article 5, holds
$\mu$ constant as far as regards the temperature because, ``in the case of small vibrations, the
terms introduced by supposing it to vary with the temperature would involve the square of the
velocity, which is supposed to be neglected.'' The reason is a bound on the velocity, and it holds
for small vibrations. Navier~\cite{src:Navier-1827} reports from the experiments he cites that the
action of the wall on the fluid diminishes as the temperature rises, and his equations hold his
constants fixed.

\textbf{A gap, stated here.} Proposition 2 puts a nonzero $\nabla T$ somewhere in every interval of
motion. Proposition 3 needs it at a point where $D\nabla T \neq 0$. Whether a flow can keep its own
temperature gradient inside the null space of its own rate of strain at every point and time is not
settled in this chapter.

\section{Condition (2) and thermal expansion}

\textbf{Proposition 4.} Let $\rho = \rho(T)$, let the mass balance and the energy balance hold, and
let $\theta = 0$. At every point where $\rho'(T) \neq 0$,
\[
\nabla\cdot(k\nabla T) = -2\mu\,|D|^2 \le 0 .
\]

\emph{Proof.} The mass balance with $\theta = 0$ is $\rho'(T)\,(\partial_t T + u\cdot\nabla T) = 0$.
Where $\rho'(T)\neq 0$ the material rate of $T$ is zero, and the energy balance leaves
$0 = \nabla\cdot(k\nabla T) + 2\mu|D|^2$. \hfill\qedsymbol

Condition (2), held together with a density that depends on temperature, therefore requires heat
conduction to remove exactly the heat that shear puts in, at every point and every moment. A fluid
that expands when heated does not do that. At a temperature where $\rho'$ vanishes the proposition
says nothing, and water has one: IAPWS R6-95 (2018)~\cite{src:IAPWS-R6-95-2018} puts its density maximum at $3.978\,^\circ$C at
atmospheric pressure. That is the only temperature at
which condition (2) and thermal expansion agree without a further condition.

\section{Where dropping the term is least safe}

Dropping $2D\nabla\mu$ claims that it is small, and $D$ is the factor that an unbounded solution
makes large.

\textbf{Proposition 5.} Let $(u,p)$ be a smooth solution of (1) and (2) on $\R^3\times[0,T^*)$ with
$u(\cdot,t)$ supported in one compact set for every $t$, with $f$ smooth and of compact support in
$\R^3\times(0,T^*)$, with $\sup_t\|u(t)\|_{L^2} < \infty$, and with
$\limsup_{t\to T^*}\|u(t)\|_{L^\infty} = \infty$. These are the hypotheses of Theorem~1.1 of
\emph{Finite time blowup for Navier--Stokes} with $T^* = 1$. Write $Q(t) = \nu\int 2|D|^2\,dx$, the
rate at which shear puts heat into the fluid. Then
\begin{enumerate}
\item $Q(t) = \nu\,\|\nabla u(t)\|_{L^2}^2$;
\item $Q$ has no finite bound on $[0,T^*)$;
\item $\int_0^{T^*} Q(t)\,dt < \infty$.
\end{enumerate}

\emph{Proof.} For the first,
$\int 2|D|^2 = \int|\nabla u|^2 + \int \partial_j u_i\,\partial_i u_j$, and integration by parts with
compact support gives $\int \partial_j u_i\,\partial_i u_j = -\int u_i\,\partial_i\theta = 0$.

For the second, the classical local theory gives, for data in $H^1$ with zero divergence and a smooth
force, a smooth solution on an interval whose length has a positive lower bound depending only on
$\|\nabla u\|_{L^2}$ at the start and on bounds of the force. If $\|\nabla u(t)\|_{L^2} \le M$ on
$[0,T^*)$, starting that theory at a time closer to $T^*$ than the bound for $M$ continues the
solution smoothly past $T^*$, uniqueness makes the continuation agree with $u$, and $\|u\|_{L^\infty}$
stays bounded near $T^*$, which contradicts the hypothesis. This step rests on the local theory,
which is not reproduced here. Leray~\cite{src:Leray-1934} states it in his Section 19: a solution
regular before $T$ that becomes irregular at $T$ has $\max|u(\cdot,t)| > A\sqrt{\nu/(T-t)}$.

For the third, a smooth solution with compact support satisfies the energy equality
\[
\tfrac12\|u(t)\|^2 + \nu\int_0^t\|\nabla u\|^2 = \tfrac12\|u(0)\|^2 + \int_0^t\!\!\int f\cdot u ,
\]
and the right side is at most
$\tfrac12\|u(0)\|^2 + \sup_s\|u(s)\|_{L^2}\int_0^{T^*}\|f(s)\|_{L^2}\,ds$. \hfill\qedsymbol

In the construction that settles statement (C), the heat that shear puts into the fluid arrives at a
rate with no bound as the singular time approaches, while its total stays finite. Equation (1) has
no channel through which that heat reaches the viscosity, and Proposition 3 says that channel is
$2\mu'(T)\,D\nabla T$, with $D$ the factor that has no bound. An argument that the dropped term is
small needs a bound on each of its factors. Of the three, $D$ is the factor equation (1) itself
determines, and it has none.

\textbf{What Proposition 5 does not show.} It does not show that the temperature gradient becomes
large in the collapsing region, that keeping the coupling stops the collapse, or that the coupled
system is regular. The workbook carries a scaling estimate on the first of these, marked as theory.

\section{What keeping a magnitude does}

The stress law can keep what equation (1) drops. A viscosity that depends on the magnitude of the
rate of strain, $\mu(|D|) = \mu_0|D|^{p-2}$ with $p > 1$, puts a vector magnitude where equation (1)
puts a constant, and $p = 2$ returns equation (1). Lienstromberg, Schiffer and Schubert~\cite{src:Lienstromberg-Schiffer-Schubert-2025} state that
for $p > 11/5$ in three dimensions their approximations converge strongly to a solution that obeys
the energy equality, that existence of solutions obeying an energy equality for such exponents
traces back to Ladyzhenskaya~\cite{src:Ladyzhenskaya-1967,src:Ladyzhenskaya-1968,src:Ladyzhenskaya-1969},
and that in the Newtonian case, $p = 2$, it is still not known whether a solution obeying an energy
equality exists even for very regular initial data.

Which dependence the stress law keeps therefore changes which statements can be proved. This
section shows that the choice decides the question. It does not show which stress law water obeys:
the dependence of water's viscosity measured in this chapter is on temperature, not on $|D|$.

\section{Putting the fluid in a box}

Statements (B) and (D) pose the problem on $\R^3/\Z^3$: the fluid in a periodic box of side one.
Statement (D) is reached in the forced construction by placing the compact supports inside one
cube and summing integer translates. Whether the box is a container the fluid is in, or one it is
put in, decides what it keeps.

Write $\hat u(\xi) = \int_{\R^3} u(x)\,e^{-2\pi i\,\xi\cdot x}\,dx$, and for a rapidly decaying smooth
$u$ on $\R^3$ let $(Pu)(x) = \sum_{k\in\Z^3} u(x+k)$, the field put in the box.

\textbf{Proposition 6.} $Pu$ is smooth and periodic, and its Fourier coefficients are
$\widehat{Pu}(m) = \hat u(m)$ for $m\in\Z^3$. The box therefore keeps $\hat u$ on the integer
lattice and nothing between its points. There are nonzero smooth rapidly decaying fields $u$ with
$\nabla\cdot u = 0$ and $Pu = 0$. If $u$ is supported inside one open unit cube, $Pu$ equals $u$
on that cube and nothing is lost.

\emph{Proof.} The coefficient is $\int_{[0,1)^3}\sum_k u(x+k)\,e^{-2\pi i\,m\cdot x}\,dx$, and since
$e^{-2\pi i\,m\cdot k} = 1$ the sum over cells unfolds to $\int_{\R^3} u(x)\,e^{-2\pi i\,m\cdot x}\,dx
= \hat u(m)$. A smooth periodic field is fixed by its coefficients, so $Pu = 0$ exactly when $\hat u$
vanishes on $\Z^3$. Take $\varphi$ smooth, nonzero, supported in the two balls of radius $1/4$
about $\pm(1/2,0,0)$, with $\varphi(-\xi) = \overline{\varphi(\xi)}$ so that the field it makes is
real. Let $A$ be the field with $\hat A = \varphi\,e_3$, and set $u = \nabla\times A$. Then
$\hat u(\xi) = 2\pi i\,\xi\times\hat A(\xi)$ is nonzero where $\varphi$ is and vanishes on $\Z^3$, and
$\nabla\cdot u = 0$. For a field supported in one open cube the translates do not overlap.
\hfill\qedsymbol

Distinct flows on $\R^3$, as many as there are such $\varphi$, all enter the box as the same flow.
The box loses nothing exactly when the fluid already fits inside one cell of it, and the forced
construction uses that case. For a fluid larger than the box, the box keeps a lattice of samples of
its transform and drops the rest.

\textbf{Proposition 7.} The rotations and reflections of $\R^3$ that fix the origin form an infinite
group. Those that carry $\Z^3$ to itself, and so act on the box, number 48.

\emph{Proof.} An orthogonal map carrying $\Z^3$ to itself sends each $e_j$ to an integer vector of
length one, which is $\pm e_k$. Such a map is a permutation of the three axes with a sign on each:
$3!\cdot 2^3 = 48$. \hfill\qedsymbol

On $\R^3$ the expansion of a field in spherical harmonics about a point commutes with every
rotation about that point, and the degree of a harmonic is then a property of the field. In the box only
48 of those rotations act. A wave pattern read in spherical harmonics on $\R^3$ is read against
every direction; in the box it is read against three axes chosen before the fluid was looked at.
This research paper does not show that spherical harmonics recover the container from the motion.
That is a claim of the thought experiment, and the workbook holds it as theory.

\section{What this establishes and what it does not}

\textbf{Established.} The passage from the balance laws to equation (1) removes the terms of
Proposition 1 and the whole energy balance. For a fluid whose viscosity depends on temperature,
the removed viscous term is zero only where the temperature gradient lies in the null space of the
rate of strain. Any viscous motion that decays at infinity makes a temperature gradient. Condition
(2), with a density that depends on temperature, forces conduction to cancel shear heating at every
point. In the class of the forced construction, the rate of shear heating has no bound at the
singular time. Putting a fluid larger than the box into it keeps its transform on a lattice and
drops the rest, and the box keeps 48 of the rotations about a point. None of the dropped terms is
zero, and none was dropped with a bound that travels
with the equation.

\textbf{Not established.} Nothing here touches the correctness of the proof of Theorem~1.1. Its
statements, the order in which its constants are chosen and its closing estimates were read; its
technical bounds were followed and not checked line by line, and nothing read bears against it. A
theorem about equation (1) is a theorem about equation (1), and
the theorem stands as one. What this chapter bounds is what such a theorem is about: an equation
from which the temperature has been removed, and not a fluid that has one. Questions (A) and (B)
remain well posed questions about equation (1), and nothing here answers either.

The reasons for the dropped terms are stated where the equation was first written.
Navier~\cite{src:Navier-1827} takes for principle that the repulsion between two molecules is
increased or diminished by a quantity proportional to the speed at which they approach or separate,
and his fluid is incompressible from the start. Stokes~\cite{src:Stokes-1845} sets the bulk term
to zero by assumption and writes that agreement with experiment must not be regarded as confirming
it; he holds $\mu$ constant in temperature for small vibrations, as above, and independent of
pressure on Du Buat's experiments in pipes. Gibbs~\cite{src:Gibbs-1884} gives the operators the
propositions above use, and the split of a field into solenoidal and irrotational parts, unique on
all of space.

The thought experiment that prompted this chapter reports a finding that a collapsing vortex in
water would vaporize near 0.7 nanometers of width. No source for that report was located. The
nearest source found, Duraiswami~\cite{src:Duraiswami-2026}, puts the first failure of the
continuum in water at cavitation near a millimeter of core radius, gives 0.3 nanometers as the
molecular scale, and gives 0.7 micrometers only as the width at which the Knudsen number in air
reaches $10^{-1}$. Nothing here rests on the report.
